\begin{figure}[t]
\centering
\begin{tikzpicture}
\begin{axis}[txaiaxis,
  xmode=log, log basis x=10,
  xmin=0.85, xmax=2400, ymin=-0.03, ymax=1.06,
  ytick={0,0.25,0.5,0.75,1},
  xlabel={Features deleted in attribution order, $k$},
  ylabel={Mean detector confidence},
  legend style={at={(0.02,0.02)},anchor=south west}]
\addplot[black!45,dashed,forget plot,line width=0.5pt,mark=none,domain=0.85:2400]
  {0.95362};
\addplot[black!45,dotted,forget plot,line width=0.5pt,mark=none,domain=0.85:2400]
  {0.5};
\node[anchor=east,font=\tiny,text=black!60,inner sep=1.5pt]
  at (axis cs:2300,0.93) {undeleted mean};
\node[anchor=west,font=\tiny,text=black!60,inner sep=1.5pt]
  at (axis cs:1.05,0.55) {decision boundary};
\addplot[color=oiblue,mark=*] coordinates {(1,0.901103) (2,0.858396) (5,0.769346) (10,0.606919) (20,0.248555) (50,0.051188) (100,0.0173434) (200,0.00707614) (500,0.00385489) (1000,0.0033043) (2000,0.00332744)};
\addlegendentry{TreeSHAP}
\addplot[color=oiorange,mark=square*] coordinates {(1,0.930712) (2,0.909883) (5,0.849614) (10,0.82172) (20,0.793908) (50,0.74693) (100,0.692126) (200,0.60632) (500,0.37843) (1000,0.15033) (2000,0.0185799)};
\addlegendentry{LIME}
\addplot[color=black,mark=diamond*] coordinates {(1,0.95362) (2,0.95362) (5,0.953351) (10,0.953426) (20,0.953354) (50,0.932046) (100,0.935278) (200,0.925253) (500,0.788027) (1000,0.523355) (2000,0.0163025)};
\addlegendentry{constant}
\addplot[color=oipurple,mark=triangle*] coordinates {(1,0.953371) (2,0.952679) (5,0.952558) (10,0.953206) (20,0.952077) (50,0.94997) (100,0.944528) (200,0.922825) (500,0.828201) (1000,0.551126) (2000,0.0465364)};
\addlegendentry{random}
\end{axis}
\end{tikzpicture}
\caption{Deletion behavior of the four explanation maps on the $200$
alerts of the instantiation of Section~\ref{sec:numsetup}: mean confidence
of the \textsf{histgb} detector after deleting the $k$ highest-attributed
features. The two gray guides mark the undeleted mean confidence
($0.9536$) and the decision boundary. Read the
curves for their \emph{ordering and shape}, not for a per-alert guarantee:
TreeSHAP crosses the boundary at $k=20$, LIME at
$k=500$, and neither control before $k=2000$, by which point almost the entire
feature vector has been deleted and the crossing says nothing about the ordering
that produced it. LIME's curve is
measured on the reduced grid its cost forced (Section~\ref{sec:numlimits}) and
is not directly comparable to the others. Single corpus, single detector, single
seed; the ordering is not asserted beyond this instantiation.}
\label{fig:deletion}
\end{figure}

\begin{figure}[t]
\centering
\begin{tikzpicture}
\begin{groupplot}[txaiaxis,
  % Two panels, not one axis with a break: the ranges are disjoint and nothing
  % here licenses reading a vertical distance across them. ``vertical sep`` is
  % wide enough that the pair does not read as a single interrupted scale.
  group style={group size=1 by 2, vertical sep=17pt,
                x descriptions at=edge bottom},
  % ``scale only axis'' so that the two heights below size the plotting boxes
  % themselves: without it pgfplots subtracts the shared x labels from the short
  % lower panel and aborts with a negative plot height. The width then excludes
  % the y label and tick labels, hence 0.82 rather than 0.97:
  % at 0.84 the picture runs 2.4pt past the column.
  scale only axis, width=0.82\columnwidth,
  xmode=log, log basis x=10,
  xtick={0.001,0.005,0.01,0.05},
  xticklabels={$0.1\%$,$0.5\%$,$1\%$,$5\%$},
  x tick label style={/pgf/number format/assume math mode=true},
  y tick label style={/pgf/number format/assume math mode=true},
  xmin=0.0008, xmax=0.065]
\nextgroupplot[height=0.40\columnwidth, ymin=0.835, ymax=1.025,
  ytick={0.85,0.9,0.95,1},
  ylabel={$B_{\Gamma,r}$: maps under audit},
  legend style={at={(0.03,0.05)},anchor=south west}]
\addplot[color=oiblue,mark=*] coordinates {(0.001,0.925359) (0.005,0.901344) (0.01,0.884201) (0.05,0.833482)};
\addlegendentry{\textsf{append\_bytes}}
\addplot[color=oigreen,mark=square*] coordinates {(0.001,0.983732) (0.005,0.958758) (0.01,0.942361) (0.05,0.875939)};
\addlegendentry{\textsf{add\_imports}}
\addplot[color=oisky,mark=triangle*] coordinates {(0.001,0.981141) (0.005,0.953819) (0.01,0.932973) (0.05,0.858819)};
\addlegendentry{\textsf{generic}}
\addplot[color=black,mark=diamond*,dashed] coordinates {(0.001,1) (0.005,1) (0.01,1) (0.05,1)};
\addlegendentry{constant control}
\nextgroupplot[height=0.16\columnwidth,
  ymin=0.5602, ymax=0.5902,
  ytick={0.5652,0.5752,0.5852}, yticklabels={$0.565$,$0.575$,$0.585$},
  ylabel={$B_{\Gamma,r}$: random},
  xlabel={Perturbation budget $r$ (fraction of the input)}]
\addplot[color=oipurple,mark=triangle*,dashed] coordinates {(0.001,0.574958) (0.005,0.575177) (0.01,0.575433) (0.05,0.575207)};
\end{groupplot}
\end{tikzpicture}
\caption{Robustness score \eqref{eq:robustness} against the adversary's budget,
each point a mean over $500$ alerts $\times$ $20$ sampled transformations. Both
panels plot the same quantity on the same horizontal axis, on \emph{separate}
vertical ranges: the random control lies far below every other map, and on one
continuous scale it flattens the three curves that actually move into the top
seventh of the panel. The panels are drawn apart rather than as one axis with a
break because no comparison of vertical distances across them is intended or
supported; read each panel's movement against its own scale, and the two
absolute levels from the tick labels. \emph{Upper}: TreeSHAP under the three
admitted transformation families (solid) and the constant control (dashed), the
controls measured under \textsf{append\_bytes}. \emph{Lower}: the random
control, on a range $0.03$ wide, so a flat line here is flat and not a range
fitted to noise. Section~\ref{sec:numrobust} reads the three movements and says
what the flat lines do and do not license; every value is computed from a maximum
over $20$ sampled members of $\Delta_z$ rather than the supremum over all of it,
and so over-estimates $B_{\Gamma,r}$. LIME
is absent: its cost admitted only the single $1\%$ point, reported in the text.
Single corpus, single detector, single seed.}
\label{fig:robustness}
\end{figure}

\begin{figure}[t]
\centering
\begin{tikzpicture}
\begin{axis}[txaiaxis,
  xmin=-0.03, xmax=1.03, ymin=-0.05, ymax=1.16,
  ytick={0,0.25,0.5,0.75,1},
  xlabel={Disclosure threshold $\tau_d$ (at $\tau_f=0.5$, $\tau_b=0.8$)},
  ylabel={Fraction of alerts admissible},
  legend style={at={(0.03,0.70)},anchor=west}]
\addplot[color=black,mark=diamond*] coordinates {(0,0) (0.1,0) (0.2,0) (0.3,0) (0.4,0) (0.5,0) (0.6,0) (0.7,0.676) (0.8,0.676) (0.9,0.676) (1,0.676)};
\addlegendentry{constant}
\addplot[color=oiblue,mark=*] coordinates {(0,0) (0.1,0) (0.2,0) (0.3,0) (0.4,0) (0.5,0.004) (0.6,0.028) (0.7,0.222) (0.8,0.312) (0.9,0.754) (1,0.758)};
\addlegendentry{TreeSHAP}
\addplot[color=oiorange,mark=square*] coordinates {(0,0) (0.1,0) (0.2,0) (0.3,0) (0.4,0.01) (0.5,0.01) (0.6,0.02) (0.7,0.02) (0.8,0.02) (0.9,0.02) (1,0.02)};
\addlegendentry{LIME}
\addplot[color=oipurple,mark=triangle*] coordinates {(0,0) (0.1,0) (0.2,0) (0.3,0) (0.4,0) (0.5,0) (0.6,0) (0.7,0) (0.8,0) (0.9,0) (1,0)};
\addlegendentry{random}
\draw[<->,black!55,line width=0.5pt]
  (axis cs:0.7,0.262) --
  (axis cs:0.7,0.636);
\node[anchor=south,font=\scriptsize,text=black!70,inner sep=2pt]
  at (axis cs:0.7,0.696)
  {$\tau_d=0.7$: $0.68$ vs $0.22$};
\end{axis}
\end{tikzpicture}
\caption{The four numeric conditions of Definition~\ref{def:admissible}, applied
to the alerts of Section~\ref{sec:numsetup} as the disclosure threshold is
relaxed with the other two thresholds held at the values used throughout
Section~\ref{sec:numerical} (the analyst role; the other three roles coincide
with it to within a percentage point on this slice). This is the measurement
behind the calibration argument: the marked span is the whole of it. At $\tau_d=0.7$ the constant control is admissible on $0.676$ of the alerts against the Shapley attribution's $0.222$---a map that explains nothing accepted on $3.0$ times as many alerts as one that does. The ordering reverses at $\tau_d=0.9$ ($0.754$ against $0.676$), so the acceptance of the degenerate map is confined to a band of the disclosure threshold rather than holding across it. The band is a property of this
instantiation, not a theorem, and it is a property of \emph{this} $F$ and \emph{this} $D$: Section~\ref{sec:controls} re-measures the comparison under three instantiations of each and finds the control ahead in 2 of the 9 combinations. What generalizes is therefore not
that the numeric conditions cannot exclude the degenerate map, but that whether
they exclude it is decided by how $F$ and $D$ are instantiated---which is why
\eqref{eq:calibration} and the contract must pin that choice down rather than
leave it to the deployment. Single corpus, single detector, single seed.}
\label{fig:admissible}
\end{figure}

\begin{figure}[t]
\centering
\begin{tikzpicture}
\begin{axis}[txaiaxis,
  xmode=log, ymode=log, log basis x=10, log basis y=10,
  xmin=1.2e-3, xmax=0.9, ymin=2.5e-5, ymax=0.9,
  xlabel={Right-hand side of \eqref{eq:bridge}},
  ylabel={Measured $\mathrm{Adv}^{\mathrm{TV}}_{E,\Gamma}$},
  legend style={at={(0.03,0.88)},anchor=west}]
\addplot[black,densely dashed,mark=none,domain=1.2e-3:0.9,samples=2,
  forget plot] {x};
\addplot[only marks,color=oipurple,mark=o,mark size=1.3pt] coordinates
  {(0.0746408,0.0106002) (0.0746408,0.00420214) (0.0746408,0.00113473) (0.0986561,0.0159597) (0.0986561,0.00758012) (0.0986561,0.00212088) (0.115799,0.0217955) (0.115799,0.0102713) (0.115799,0.00290447) (0.166518,0.0364056) (0.166518,0.0170076) (0.166518,0.00483306) (0.0162684,0.00102914) (0.0162684,0.000515547) (0.0162684,0.000133099) (0.0412417,0.00353031) (0.0412417,0.00117087) (0.0412417,0.000288321) (0.0576387,0.00717336) (0.0576387,0.00232291) (0.0576387,0.000577451) (0.124061,0.029586) (0.124061,0.00973439) (0.124061,0.00246987) (0.0188591,0.000300755) (0.0188591,0.000165496) (0.0188591,5.83294e-05) (0.0461812,0.00084275) (0.0461812,0.000315813) (0.0461812,9.1636e-05) (0.067027,0.001415) (0.067027,0.000596561) (0.067027,0.00018254) (0.141181,0.00539505) (0.141181,0.00297435) (0.141181,0.000970378) (0.402217,0.000947362) (0.402217,0.00030968) (0.402217,8.75694e-05) (0.425042,0.00105975) (0.425042,0.000460034) (0.425042,0.000126828) (0.424823,0.0012228) (0.424823,0.000482225) (0.424823,0.0001249) (0.424567,0.00119147) (0.424567,0.000434154) (0.424567,0.000117384) (0.424793,0.00130513) (0.424793,0.000429246) (0.424793,0.000114367) (0.424545,0.00110566) (0.424545,0.000459954) (0.424545,0.00012363) (0.42483,0.00111283) (0.42483,0.000454543) (0.42483,0.00012215) (0.424868,0.00137604) (0.424868,0.00050697) (0.424868,0.000154052) (0.424822,0.00111077) (0.424822,0.000421214) (0.424822,0.000125858) (0.424613,0.00106067) (0.424613,0.000387588) (0.424613,0.000102255) (0.424929,0.001016) (0.424929,0.000421201) (0.424929,0.000122919) (0.424849,0.000965683) (0.424849,0.000359672) (0.424849,0.000103406) (0.424834,0.00145566) (0.424834,0.000523977) (0.424834,0.000141861)};
\addlegendentry{generic constant $L_\pi=1$}
\addplot[only marks,color=oiblue,mark=x,mark size=1.9pt] coordinates
  {(0.0733174,0.0106002) (0.0451951,0.00420214) (0.0143882,0.00113473) (0.0969069,0.0159597) (0.0597363,0.00758012) (0.0190175,0.00212088) (0.113746,0.0217955) (0.0701163,0.0102713) (0.022322,0.00290447) (0.163566,0.0364056) (0.100827,0.0170076) (0.032099,0.00483306) (0.0159799,0.00102914) (0.00985052,0.000515547) (0.00313598,0.000133099) (0.0405105,0.00353031) (0.0249719,0.00117087) (0.00794998,0.000288321) (0.0566168,0.00717336) (0.0349003,0.00232291) (0.0111108,0.000577451) (0.121862,0.029586) (0.0751192,0.00973439) (0.0239147,0.00246987) (0.0185247,0.000300755) (0.0114192,0.000165496) (0.00363539,5.83294e-05) (0.0453624,0.00084275) (0.0279627,0.000315813) (0.00890214,9.1636e-05) (0.0658386,0.001415) (0.0405849,0.000596561) (0.0129205,0.00018254) (0.138678,0.00539505) (0.0854854,0.00297435) (0.0272149,0.000970378) (0.395086,0.000947362) (0.243543,0.00030968) (0.0775337,8.75694e-05) (0.417505,0.00105975) (0.257363,0.000460034) (0.0819334,0.000126828) (0.417291,0.0012228) (0.257231,0.000482225) (0.0818913,0.0001249) (0.417039,0.00119147) (0.257076,0.000434154) (0.0818419,0.000117384) (0.417262,0.00130513) (0.257213,0.000429246) (0.0818856,0.000114367) (0.417018,0.00110566) (0.257062,0.000459954) (0.0818377,0.00012363) (0.417298,0.00111283) (0.257235,0.000454543) (0.0818927,0.00012215) (0.417335,0.00137604) (0.257258,0.00050697) (0.0819,0.000154052) (0.41729,0.00111077) (0.25723,0.000421214) (0.0818912,0.000125858) (0.417084,0.00106067) (0.257103,0.000387588) (0.0818508,0.000102255) (0.417395,0.001016) (0.257295,0.000421201) (0.0819116,0.000122919) (0.417316,0.000965683) (0.257246,0.000359672) (0.0818962,0.000103406) (0.417301,0.00145566) (0.257237,0.000523977) (0.0818933,0.000141861)};
\addlegendentry{Dobrushin coefficient $\delta$}
\end{axis}
\end{tikzpicture}
\caption{Theorem~\ref{thm:bridge} on the instantiated kernel. Each of the
$75$ plotted configurations lies below the diagonal $y=x$ (dashed), so the
inequality is not violated, and each lies far below it, so the bound is
loose. Replacing the generic constant $L_\pi=1$ by the ergodic
coefficient \eqref{eq:dobrushin} moves each point left without moving it above
the diagonal, which is the sharpening claimed in the text. The remaining $36$ of the $111$ recomputed
configurations are the constant control's, omitted because both sides are
exactly zero and so have no position on logarithmic axes; they are also the only
configurations in which the bound is attained. Points are configurations, not alerts, and
holding on the configurations of one instantiation is not a claim that the
bound is tight or loose in general.}
\label{fig:bridge}
\end{figure}

\begin{table}[t]
\caption{Strong Stackelberg equilibrium of the $15\times13$ release game
of Section~\ref{sec:numgame}, over the two weights the utilities of
Definition~\ref{def:game} leave free: $\eta$ on evidence failure and $\rho$ on
disclosure. $k^\star$ is the equilibrium release width on the lattice
$\{5,10,20,50,100\}$ and $U_D$ the defender's equilibrium value. The width is
nondecreasing in $\eta$ and nonincreasing in $\rho$: an evidence obligation is
the only term that pays for release, and at $\eta=0$ the equilibrium releases
the least the lattice allows at every $\rho$. The equilibrium is pure at all
$18$ settings---the gain from randomizing does not
exceed $1.2\times10^{-16}$---and the attacker plays pure observation
($M_4$) at every one of them.}
\label{tab:numgame}
\centering
\footnotesize
\begin{tabular}{@{}lcccccc@{}}
\toprule
& \multicolumn{2}{c}{$\eta=0$} & \multicolumn{2}{c}{$\eta=1$} & \multicolumn{2}{c}{$\eta=4$}\\
$\rho$ & $k^\star$ & $U_D$ & $k^\star$ & $U_D$ & $k^\star$ & $U_D$\\
\midrule
$0$ & $5$ & $-0.339$ & $100$ & $-0.412$ & $100$ & $-0.484$\\
$0.25$ & $5$ & $-0.368$ & $100$ & $-0.604$ & $100$ & $-0.676$\\
$0.5$ & $5$ & $-0.396$ & $100$ & $-0.796$ & $100$ & $-0.868$\\
$1$ & $5$ & $-0.453$ & $20$ & $-1.076$ & $100$ & $-1.252$\\
$2$ & $5$ & $-0.566$ & $10$ & $-1.440$ & $100$ & $-2.019$\\
$4$ & $5$ & $-0.792$ & $5$ & $-1.695$ & $20$ & $-3.181$\\
\bottomrule
\end{tabular}
\end{table}
